\chapter{Myocarditis}

\medskip
\noindent\textit{\textbf{Under review.} This chapter is an AI-assisted educational draft; it is not a substitute for professional medical advice.}

\section*{Definition and Overview}
inflammation of the heart muscle, usually caused by infection or an immune-mediated reaction. Individual assessment is important because symptoms and treatment vary with the cause.

\section*{Clinical Features}
Chest pain, breathlessness, palpitations, fatigue, fainting, or reduced exercise tolerance.

\section*{When to Seek Medical Care}
Seek medical review for new, persistent, worsening, or function-limiting symptoms. Seek emergency help for severe breathing difficulty, collapse, sudden neurological change, severe chest or abdominal pain, or any rapidly worsening illness.

\section*{Diagnosis and Investigations}
Electrocardiogram, troponin and other blood tests, echocardiography, cardiac MRI, and sometimes coronary assessment or biopsy.

\section*{Management}
Rest and avoidance of strenuous exercise, treatment of heart failure or arrhythmia, and specialist management of the underlying cause. Do not resume sport until medically cleared.

\section*{Complications}
Possible complications depend on the cause and may include progression of the condition, effects on other organs, treatment adverse effects, and reduced quality of life. Early assessment can reduce preventable harm.

\section*{Prognosis and Outcomes}
Outcome depends on the underlying cause, severity, complications, and response to treatment. Discuss individual prognosis with the treating clinician.

\section*{Prevention and Self-Care}
Follow the treatment plan, keep relevant appointments, avoid known triggers, protect affected areas from injury, and seek help early if symptoms worsen. Do not start or stop prescription treatment without clinical advice.

\clearpage
